\documentclass{article}
\usepackage[T1]{fontenc}
\usepackage{lmodern}
\usepackage{graphicx}
\usepackage{amsmath,amssymb}
\usepackage[a4paper,margin=2cm]{geometry}
\usepackage[absolute,overlay]{textpos}
\setlength{\TPHorizModule}{1mm}
\setlength{\TPVertModule}{1mm}
\usepackage[indonesian]{babel}
\usepackage{hyperref}
\setlength{\emergencystretch}{2em}
% unit: Halaman 39

\begin{document}

% === Halaman 39 (python/native) ===

• Misalkan kita mengetahui pasangan potongan plaintext (p) dan cipherteks (c) 

yang berkoresponden. 

• Lakukan serangan brute force (atau exhaustive key search attack) dengan 

mencoba semua kemungkinan kunci (k):

for(k=0;k<256;k++) \{

 if(DES(k,p)==c) \{


printf(“Key is \%x\textbackslash{}n”, k);


break;

 \}

\}

• Kunci diharapkan ditemukan sebelum 255 iterasi.

39

\end{document}
